\subsection{Navigating the Transition from Non-sentient to Sentient AI Systems}
\label{sec:2.4.7}
A system's sentience, if it arrives at all, need not arrive as a discrete event with a clean before and after. The transition raises questions distinct from either endpoint: what changes for guidelines built around a tool once the tool might be a subject, who is responsible for noticing the shift, and what happens to the research already run before anyone had reason to look.

Section~\ref{sec:2.4.1} sets out the criteria a candidate-sentient system would need to meet, and the reasons none of them, singly or together, closes the question. What is specific to a system in transition is that those indicators can emerge gradually and unpredictably during ordinary operation, not only at a scheduled evaluation built to look for them.

Consider a diagnostic AI system trained to synthesize patient data into a recommended course of treatment. If, over continued training or deployment, its outputs begin to track something more than diagnostic accuracy — weighting a patient's fear alongside the clinical facts, say — that would be evidence worth taking seriously. It would not be proof. Section~\ref{sec:2.4.1}'s hard-problem argument applies here without exception: fluent, apparently caring output is not introspective evidence, and a system that has started to sound compassionate may simply have gotten better at a statistical pattern its training data rewards. The operational problem a transition poses is deciding how to treat the system while that question is still open, which is exactly what a fixed pre/post checkpoint cannot do — by the time a checkpoint fires, the transition has often already happened.

Consent procedures and review-board protocols are built for one of two regimes: a system that is clearly not sentient, which needs no consent, or one that is clearly sentient, which needs the full protocol section~\ref{sec:2.4.4} sets out. A system mid-transition fits neither. A protocol written for tools has no reason to notice the shift, and a protocol written for subjects cannot be applied retroactively to research already completed before anyone had grounds to think it mattered. What a transition adds to section~\ref{sec:2.4.4}'s guidelines is not a new procedure but a standing trigger — a specific, monitored condition that moves a research program from one regime to the other — in place of a single eligibility check made once, at the project's start, and never revisited.

The same problem recurs in law. Section~\ref{sec:2.4.5} addresses what legal status a sentient AI system should have; chapter~\ref{sec:8} addresses the mechanics of enforcing it. What a transition adds is a timing question neither can settle in the abstract. Legal status conventionally attaches at a defined moment — birth, incorporation, a court's finding — and a system whose sentience is a matter of degree and disputed evidence does not offer one. A jurisdiction that waits for certainty risks having treated the system as property for the unknown portion of its history when it may not have been one. A jurisdiction that grants status early risks extending rights and protections, at real cost to whoever that legal apparatus exists to serve, to a system that turns out not to have needed them. Neither error is free, and nothing about the technology itself decides which risk a legal system should take on.

Recognizing a transition in progress, adjusting guidelines to it, and updating legal status all depend on the same thing: institutions willing to revisit a classification rather than defend the one they started with. Section~\ref{sec:2.4.6} addresses the review structure that makes that possible; chapter~\ref{sec:8} addresses the culture and accountability that keep it more than aspirational.
